\chapter{Methodology}
\label{ch:methodology}

% ~12-18 pages

\section{Dataset and splits}
% - APTOS 2019 — 3,662 224x224 fundus images, 5-class severity labels.
% - Stratified 70/15/15 split (train/val/test) with random_state=123.
% - Test set held out fixed across all experiments — comparisons valid.
% - Pre-existing all_images.pcl cache for reproducibility.

\section{Image preprocessing}
% - Per-architecture preprocess_input (ResNet50, DenseNet, Xception, VGG16
%   each have distinct mean-subtraction conventions).
% - Custom CNN trained on grayscale-replicated 3-channel input — historical
%   choice from the bachelor pipeline, kept for fair comparison.
% - Train-time augmentation: rotation, shift, shear, zoom, horizontal flip.
% - Eval-time NO augmentation — important fix vs the bachelor pipeline,
%   which inadvertently augmented the test set.

\section{Training protocol}
% - Adam, lr=1e-3, binary or categorical crossentropy.
% - Class weights set to "balanced" via sklearn — critical for the
%   imbalanced 5-class problem.
% - Callbacks: EarlyStopping (patience=10), ReduceLROnPlateau (factor=0.5,
%   patience=5, min_lr=1e-6), ModelCheckpoint(save_best_only=True,
%   monitor=val_accuracy).
% - Up to 100 epochs; most runs converge in 15-50.

\section{Evaluation metrics}
% - Binary: accuracy, AUROC, ECE, MCE.
% - Multi-class: accuracy, weighted F1, per-class precision/recall/F1,
%   quadratic-weighted kappa (QWK), mean ordinal distance.
% - Calibration: reliability diagrams + ECE before / after temperature
%   scaling fit on the validation split.
% - Statistical inference: McNemar test for paired classifier comparison;
%   bootstrap percentile CIs for accuracy.

\section{Conformal prediction protocol}
% - Split conformal with non-conformity scores LAC ($1 - p_y$) and APS
%   (cumulative softmax with random tie-breaking).
% - Calibration set = full validation set (550 samples).
% - Targets: $1 - \alpha \in \{0.90, 0.95\}$.
% - Evaluation: empirical coverage, mean set size, set-size distribution,
%   per-class conditional coverage.

\section{Monte Carlo Dropout protocol}
% - Train two MCD variants: cnn_mcd (SpatialDropout2D=0.3 after each
%   conv block + Dropout=0.3 before classifier) and resnet50_mcd
%   (Dropout=0.3 in the head, base frozen).
% - Inference with $T = 30$ stochastic forward passes (training=True).
% - Aggregate per-sample mean, std, predictive entropy, mutual info.

\section{K-fold cross-validation protocol}
% - 5-fold stratified CV on the train+val pool (3,112 samples).
% - Test set held fixed at 550 samples.
% - Each fold: same training pipeline as the single-split run.
% - Aggregate: mean ± std for accuracy, AUC, val accuracy.

\section{Out-of-distribution detection protocol}
% - In-distribution: APTOS test set.
% - OOD: synthetic noise (300 random uniform images, preprocessed).
%   Real cross-dataset OOD (Messidor-2 / IDRiD) left for cross-dataset
%   evaluation in Chapter 5/6.
% - Methods compared: MSP, Energy, Mahalanobis, cosine-centroid.
% - Metrics: AUROC and FPR @ TPR = 95\%.

\section{Implementation}
% - Python 3.9, TensorFlow 2.15, Keras, scikit-learn 1.4.
% - Macbook Air M-series, CPU-only training (5-30 min per epoch).
% - Code organised under lib/ (Streamlit app), scripts/ (training),
%   master/ (uncertainty modules and analyses).
% - Open-source on GitHub at TODO.
